\subsection*{A cry gathered from within betrayal\enspace\properrefs{Int.}}

The Introit is a cento rather than a continuous extract. Its printed marginal
reference names Psalm 54:17, 18, 20, and 23, followed by verse 2 as the psalm
verse. The sung wording also preserves the clause \latin{ab his qui
appropínquant mihi}, ``from those who draw near to me,'' from verse 19. The
Missal's printed citation therefore names fewer verses than its text draws on:
the sung wording reaches beyond the verses printed in the margin.

That small textual distinction recovers the pressure inside the chant. The
complete psalm moves from urgent petition through fear and the desire to flee,
violence and deceit in the city, and betrayal by an intimate companion who had
walked in God's house. Only then come the cry to God, prayer at evening,
morning, and noon, redemption from those who draw near, judgment by the One
who abides before the ages, and the command to cast care upon the Lord. The
selected confidence does not pretend that treachery or danger is imaginary;
it transfers the burden's final custody to God while the danger remains.

Augustine's \emph{Enarratio in Psalmum 54} follows this movement rather than
extracting a placid maxim from it. Sections 16--19 treat the cry, the three
hours, deliverance amid those nearby, and the Eternal's humbling judgment.
His reading of evening, morning, and noon through Christ's Passion,
Resurrection, and Ascension is ecclesial and Christological; it rests on the
lament's human horizon and expounds the psalm, not the Roman chant. It also keeps the phrase ``those who draw near'' attached to the
painful nearness of divided former companions rather than turning it into a
license to nominate present-day enemies.

At §23 Augustine directly treats \latin{Iacta cogitátum tuum in Dómino}. His
pastoral force is dependence: the hearer casts self and care upon the Lord
rather than putting a creature in the Lord's place, and the Lord nourishes the
one not yet able to grasp everything. Section 24 begins the following verse's
descent into the pit; it supplies adjacent context but is not the direct locus
for the appointed command. Casting care therefore means entrusted dependence
on the sustaining Creator, not denial of danger, abandonment of duty, or
passivity toward a burden one can responsibly address.

The opening psalm verse returns the whole entrance chant to petition:
\latin{Exáudi, Deus, oratiónem meam}. The liturgical speaker does not first
announce mastery of fear and then approach God. The speaker asks to be heard,
rehearses the prayer that has been heard, confesses divine judgment, and
receives the imperative to entrust care. Within this Introit, confidence is
the form taken by persevering prayer under pressure.
